% ---------------------------------------------------------------------------
% Preamble for the Cassini ISS F Ring Mosaics User Guide
% Converted from users-guide-draft-phase2-6.docx
% ---------------------------------------------------------------------------

\usepackage[utf8]{inputenc}
\usepackage[T1]{fontenc}
\usepackage{textcomp}

% Base fonts: Times body (Word: Times New Roman/Minion), Helvetica for
% headings (Word: Calibri), Courier for code (Word: Cascadia Code).
\usepackage{mathptmx}
\usepackage[scaled=0.92]{helvet}
\usepackage{courier}

% Page geometry matches the Word original (US Letter,
% margins: left/right/top 0.75 in, bottom 1.25 in).
\usepackage[letterpaper,left=0.75in,right=0.75in,top=0.75in,bottom=1.25in,
            footskip=0.5in]{geometry}

\usepackage{graphicx}
\usepackage{longtable}
\usepackage{array}
\usepackage{amsmath}
\usepackage{amssymb}
\usepackage{fancyvrb}
\usepackage{fvextra}     % adds automatic line breaking to verbatim blocks
\usepackage{float}
\usepackage{enumitem}
\usepackage{titlesec}
\usepackage{caption}
\usepackage{fancyhdr}
\usepackage{pmboxdraw}      % box-drawing characters in directory listings
\usepackage{newunicodechar}
\usepackage{lscape}         % Section 7 is landscape, as in the Word original
\usepackage{lastpage}

% Rotate the physical PDF pages of landscape sections (equivalent to
% pdflscape, but implemented directly for maximum portability).
%
% lscape rotates the body of the page but not the running foot, so with the
% default page style the "Page N of M" footer ends up reading up the left edge
% of the rotated page.  The landscape page style below instead draws the
% footer rotated, out in the margin beyond the right-hand edge of the (still
% unrotated) text block and centred on the page vertically, which is where the
% bottom centre of the page lands once the physical page is turned.  The two
% offsets are measured from the page, not from the text block, so they hold
% for any page in the landscape section; \lsfootwidth fixes the length of the
% footer box so that its centring does not shift with the page number.
\newlength{\lsfootraise}   \setlength{\lsfootraise}{290pt}
\newlength{\lsfootshift}   \setlength{\lsfootshift}{20pt}
\newlength{\lsfootwidth}   \setlength{\lsfootwidth}{1.5in}
\fancypagestyle{landscapepage}{%
  \fancyhf{}%
  \renewcommand{\headrulewidth}{0pt}%
  \fancyfoot[R]{\makebox[0pt][l]{\hspace{\lsfootshift}%
    \raisebox{\lsfootraise}[0pt][0pt]{\rotatebox{90}{%
      \makebox[\lsfootwidth][c]{\small Page \thepage\ of \pageref{LastPage}}}}}}%
}

\ifdefined\pdfpageattr
  \newenvironment{rotatedlandscape}
    {\clearpage\global\pdfpageattr{/Rotate 90}\begin{landscape}%
     \pagestyle{landscapepage}\thispagestyle{landscapepage}}
    {\end{landscape}\clearpage\global\pdfpageattr{}\pagestyle{fancy}}
\else
  \newenvironment{rotatedlandscape}
    {\begin{landscape}\pagestyle{landscapepage}\thispagestyle{landscapepage}}
    {\end{landscape}\pagestyle{fancy}}
\fi

% --- characters used in the Word source ------------------------------------
\newunicodechar{°}{\textdegree}
\newunicodechar{−}{\ensuremath{-}}
\newunicodechar{×}{\ensuremath{\times}}
\newunicodechar{Δ}{\ensuremath{\Delta}}
\newunicodechar{ϖ}{\ensuremath{\varpi}}
\newunicodechar{Ω}{\ensuremath{\Omega}}
\newunicodechar{λ}{\ensuremath{\lambda}}
\newunicodechar{θ}{\ensuremath{\theta}}
\newunicodechar{≈}{\ensuremath{\approx}}

% --- paragraph layout: 0.25 in first-line indent, no inter-paragraph skip,
% as in the Word original
\setlength{\parindent}{0.25in}
\setlength{\parskip}{0pt}
\emergencystretch=1em

% --- headings: numbered to 4 levels (Word Heading 1-4); TOC to 4 levels ----
\setcounter{secnumdepth}{4}
\setcounter{tocdepth}{4}

% The article class leaves a wide gap between a number and its title in the
% table of contents (1.5em for a section, growing to 4.1em for a paragraph).
% Set each gap to just clear the widest number at that level, which is the
% same effect as the 0.6em after the numbers of the headings themselves, and
% carry the narrower widths into the indents so each level still starts where
% the level above it ends.
\makeatletter
\renewcommand*\l@section[2]{%
  \ifnum \c@tocdepth >\z@
    \addpenalty\@secpenalty
    \addvspace{1.0em \@plus\p@}%
    \setlength\@tempdima{1.0em}%
    \begingroup
      \parindent \z@ \rightskip \@pnumwidth
      \parfillskip -\@pnumwidth
      \leavevmode \bfseries
      \advance\leftskip\@tempdima
      \hskip -\leftskip
      #1\nobreak\hfil \nobreak\hb@xt@\@pnumwidth{\hss #2}\par
    \endgroup
  \fi}
\renewcommand*\l@subsection{\@dottedtocline{2}{1.0em}{1.7em}}
\renewcommand*\l@subsubsection{\@dottedtocline{3}{2.7em}{2.5em}}
\renewcommand*\l@paragraph{\@dottedtocline{4}{5.2em}{3.3em}}
\makeatother

% Sizes/weights follow the Word styles:
%   Heading 1: 16 pt bold;  Heading 2: 13 pt bold;
%   Heading 3 and 4: 12 pt bold italic
\titleformat{\section}
  {\headingfont\fontsize{16pt}{19pt}\selectfont}{\thesection}{0.6em}{}
\titlespacing*{\section}{0pt}{6pt}{12pt}
\titleformat{\subsection}
  {\headingfont\fontsize{13pt}{16pt}\selectfont}{\thesubsection}{0.6em}{}
\titlespacing*{\subsection}{0pt}{12pt}{6pt}
\titleformat{\subsubsection}
  {\headingfont\itshape\fontsize{12pt}{14pt}\selectfont}{\thesubsubsection}{0.6em}{}
\titlespacing*{\subsubsection}{0pt}{12pt}{6pt}
\titleformat{\paragraph}[hang]
  {\headingfont\itshape\fontsize{12pt}{14pt}\selectfont}{\theparagraph}{0.6em}{}
\titlespacing*{\paragraph}{0pt}{6pt}{6pt}

% --- captions: "Figure N: ..." in italics, as in the Word Caption style ----
\captionsetup{labelsep=colon,font={it},skip=6pt}

% Every table in this guide is a longtable with its caption above it; leave
% clear space between the preceding text and that caption.
\setlength{\LTpre}{18pt}
\setlength{\LTpost}{12pt}

% --- lists: compact spacing similar to Word --------------------------------
\setlist{itemsep=2pt,topsep=4pt,parsep=0pt}

% --- floats: keep figures close to the text that references them -----------
% (the defaults defer a 1.5-3 in figure by several pages in this document)
\setcounter{topnumber}{3}
\setcounter{bottomnumber}{2}
\setcounter{totalnumber}{4}
\renewcommand{\topfraction}{0.92}
\renewcommand{\bottomfraction}{0.75}
\renewcommand{\textfraction}{0.07}
\renewcommand{\floatpagefraction}{0.75}

% --- tables ----------------------------------------------------------------
\setlength{\LTpre}{8pt}
\setlength{\LTpost}{8pt}
\renewcommand{\arraystretch}{1.15}

% --- footer: "Page N of M" centered, as in the Word original ---------------
\pagestyle{fancy}
\fancyhf{}
\fancyfoot[C]{\small Page \thepage\ of \pageref{LastPage}}
\renewcommand{\headrulewidth}{0pt}
\fancypagestyle{plain}{\fancyhf{}%
  \fancyfoot[C]{\small Page \thepage\ of \pageref{LastPage}}}

% --- hyperlinks ------------------------------------------------------------
% hyperref last; URLs blue as in Word, internal cross-references black.
% The pdfa option sets the Print flag on every link annotation. PDF/A-1
% forbids non-printing annotations, so without it the PDF/A conversion in the
% Makefile silently discards all 251 links in the document.
\usepackage[pdfa,
            colorlinks=true,urlcolor=blue,linkcolor=black,citecolor=black,
            bookmarksdepth=4,
            pdftitle={Cassini ISS F Ring Mosaics User Guide},
            pdfauthor={Robert S. French}]{hyperref}

% allow URLs to break more freely and stretch slightly
\Urlmuskip=0mu plus 2mu
% \allowbreak and \code appear in heading code tokens; in PDF bookmarks drop
% the former and keep only the text of the latter
\newcommand{\pdfidentity}[1]{#1}
\pdfstringdefDisableCommands{\let\allowbreak\relax \let\code\pdfidentity}

% long URLs (e.g. in the References section) may break after any character
\expandafter\def\expandafter\UrlBreaks\expandafter{\UrlBreaks%
  \do\a\do\b\do\c\do\d\do\e\do\f\do\g\do\h\do\i\do\j\do\k\do\l\do\m%
  \do\n\do\o\do\p\do\q\do\r\do\s\do\t\do\u\do\v\do\w\do\x\do\y\do\z%
  \do\0\do\1\do\2\do\3\do\4\do\5\do\6\do\7\do\8\do\9\do\-\do\_}
